% ================================================================
% ==== FILE: content/m_spaces/m9.tex
% ================================================================

% ==============================
%  Ontology of Continua — Core
%  Meta-Space: M9
%  FULL MODULE — FINAL VERSION
% ==============================

\subsubsection{\texorpdfstring{$M_9$}{M_9} Übersicht}
\label{sec:m9-overview}

$M_9$ ist der Metaraum, der die Existenz von ermöglicht
theoretische Kontinua ($K_9$).
Während $M_8$ das Materielle, Symbolische und Institutionelle bereitstellt
Als Substrat der Zivilisationen stellt $M_9$ die strukturellen Bedingungen vor
für die Entstehung, Koexistenz und Transformation von
wissenschaftliche, mathematische und konzeptionelle Modelle.

In $M_9$ sind die Grundeinheiten nicht Agenten, Infrastrukturen oder
Institutionen, sondern \emph{Modelle}, dargestellt als strukturierte Objekte
mit internen Logiken, Einschränkungen und Transformationsregeln.

Somit ist $M_9$ der minimale Raum, in dem \emph{Theorie} ein zulässiges Kontinuum wird.

% ---------------------------------------------------------------
\subsubsection*{1. Zulässiger Konfigurationsraum \texorpdfstring{$\Omega(M_9)$}{\Omega(M_9)}}

\[
\Omega(M_9)=
\left\{
(\mathcal{M},\ \mathrm{Mor},\ A_{\mathrm{Modell}},\
\Theta_{\mathrm{coh}},\ C_{\mathrm{Theorie}},\ \Sigma_9)
\mid \text{Kohärenzbedingungen gelten}
\right\}.
\]

Komponenten:

\begin{itemize}
    \item $\mathcal{M}$ – die Menge der zulässigen Modelle
        (mathematisch, physikalisch, chemisch, biologisch, kognitiv, sozial),
    \item $\mathrm{Mor}$ – Morphismen zwischen Modellen
        (Interpretationen, Reduktionen, Äquivalenzen, Funktoren),
    \item $A_{\mathrm{model}}$ – Achsen der theoretischen Variation,
    \item $\Theta_{\mathrm{coh}}$ – Kohärenzschwellen,
    \item $C_{\mathrm{theory}}$ – Zyklen der theoretischen Verfeinerung,
    \item $\Sigma_9$ – globale strukturelle Einschränkungen des Theorieraums.
\end{itemize}

Voraussetzung für die Zulassung sind:

\begin{enumerate}
    \item Existenz stabiler Modelle: $\mathcal{M}\neq\emptyset$,
    \item Existenz sinnvoller Morphismen zwischen Modellen,
    \item nicht-triviale theoretische Achsen,
    \item Kohärenz: $d_{\mathrm{coh}}(M_i, M_j)<\Theta_{\mathrm{coh}}$ für zulässige Paare,
    \item Möglichkeit der theoretischen Entwicklung unter dem Operator $E$.
\end{enumerate}

% ---------------------------------------------------------------
\subsubsection*{2. Achsen \texorpdfstring{$A(M_9)$}{A(M_9)}}

\[
A(M_9)=
\{
A_{\mathrm{Syntax}},\
A_{\mathrm{Semantik}},\
A_{\mathrm{Abstraktion}},\
A_{\mathrm{Idealisierung}},\
A_{\mathrm{Näherung}},\
A_{\mathrm{Domäne}},\
A_{\mathrm{Logik}}
\}.
\]

Beschreibungen:

\begin{itemize}
    \item $A_{\mathrm{syntax}}$ – formale Sprachen, symbolische Strukturen,
    \item $A_{\mathrm{semantics}}$ — Interpretationsräume,
    \item $A_{\mathrm{Abstraction}}$ – Grade der Allgemeinheit,
    \item $A_{\mathrm{idealization}}$ — erlaubte Vereinfachungen,
    \item $A_{\mathrm{Approximation}}$ — Fehler- und Konvergenzstrukturen,
    \item $A_{\mathrm{domain}}$ – domänenspezifische Achsen (Physik, Chemie, Biologie…),
    \item $A_{\mathrm{logic}}$ – logische Regeln, Inferenzsysteme.
\end{itemize}

Eine notwendige Voraussetzung für modelltheoretisches Leben:

\[
\dim(A_{\mathrm{Syntax}})\ge 1,
\quad
\dim(A_{\mathrm{Semantik}})\ge 1.
\]

% ---------------------------------------------------------------
\subsubsection*{3. Potentiale \texorpdfstring{$P(M_9)$}{P(M_9)}}

\[
P(M_9)=
(F_{\mathrm{Expressivität}},\
F_{\mathrm{Konsistenz}},\
F_{\mathrm{Kohärenz}},\
F_{\mathrm{Verallgemeinerung}},\
F_{\mathrm{prädiktiv}},\
F_{\mathrm{Kompression}}).
\]

Interpretation:

\begin{itemize}
    \item $F_{\mathrm{Expressivität}}$ – Potenzial, das reichhaltige theoretische Beschreibungen ermöglicht,
    \item $F_{\mathrm{Konsistenz}}$ – Potenzial zur Gewährleistung der logischen Stabilität,
    \item $F_{\mathrm{Kohärenz}}$ – potenzielle Modelle, die gegenseitig interpretierbar bleiben,
    \item $F_{\mathrm{Generalization}}$ – Kapazität für Abstraktion und Domänenübertragung,
    \item $F_{\mathrm{predictive}}$ – Genauigkeit und empirische Angemessenheit,
    \item $F_{\mathrm{Kompression}}$ – Potenzial für minimal ausreichende Darstellungen.
\end{itemize}

Bedingungen für die Existenz von $K_9$:

\[
F_{\mathrm{Kohärenz}}>\Theta_{\mathrm{Fragmentierung}},
\qquad
F_{\mathrm{Konsistenz}}>\Theta_{\mathrm{Paradoxon}}.
\]

% ---------------------------------------------------------------
\subsubsection*{4. Schwellenwerte \texorpdfstring{$\Theta(M_9)$}{\Theta(M_9)}}

Wichtige Kohärenzschwellen:

\begin{itemize}
    \item $\Theta_{\mathrm{coh}}$ – maximal zulässige Divergenz für die Koexistenz von Modellen,
    \item $\Theta_{\mathrm{Konsistenz}}$ – Schwelle der internen logischen Stabilität,
    \item $\Theta_{\mathrm{predictive}}$ – minimale Vorhersagegenauigkeit,
    \item $\Theta_{\mathrm{paradox}}$ – Grenze, an der Inkonsistenz das Kontinuum zusammenbricht,
    \item $\Theta_{\mathrm{fragmentation}}$ — Schwelle für theoretischen Zerfall.
\end{itemize}

Das Überschreiten von $\Theta_{\mathrm{paradox}}$ erzwingt den Zusammenbruch von $\Omega(K_9)$,
analog zum institutionellen Zusammenbruch in $K_8$.

% ---------------------------------------------------------------
\subsubsection*{5. Flüsse \texorpdfstring{$J(M_9)$}{J(M_9)}}

\[
J(M_9)=
(J_{\mathrm{Inferenz}},\
J_{\mathrm{Interpretation}},\
J_{\mathrm{translation}},\
J_{\mathrm{Abstraktion}},\
J_{\mathrm{Spezialisierung}},\
\nabla_i T_9).
\]

Interpretation:

\begin{itemize}
    \item $J_{\mathrm{inference}}$ – Fluss von Schlussfolgerungen und Beweisen,
    \item $J_{\mathrm{interpretation}}$ — Neuinterpretation zwischen Modellen,
    \item $J_{\mathrm{translation}}$ – Übersetzung über formale Sprachen hinweg,
    \item $J_{\mathrm{Abstraction}}$ – Fluss in Richtung höherer Allgemeinheit,
    \item $J_{\mathrm{Spezialisierung}}$ – domänenspezifische Ableitungen,
    \item $\nabla_i T_9$ – Gradienten der theoretischen Spannung.
\end{itemize}

Stabilitätsbedingung:

\[
\oint J_{\mathrm{inference}}\,dt\ approx 0
\quad\Rightarrow\quad
\text{Modell bleibt über Zyklen hinweg logisch stabil}.
\]

% ---------------------------------------------------------------
\subsubsection*{6. Zyklen \texorpdfstring{$C(M_9)$}{C(M_9)}}

Grundlegende Zyklen:

\begin{itemize}
    \item $C_{\mathrm{Theorie}}$
        — Hypothese $\to$ Ableitung $\to$ Bewertung $\to$ Revision,
    \item $C_{\mathrm{abstraction}}$
        — konkretes $\to$ abstraktes $\to$ verallgemeinertes $\to$ instanziiert,
    \item $C_{\mathrm{Interpretation}}$
        — Modell $\to$ Mapping $\to$ Neuinterpretation,
    \item $C_{\mathrm{Kohärenz}}$
        — lokale Anpassungen $\bis$ globale Konsistenz $\bis$ Refit-Zyklen,
    \item $C_{\mathrm{Vereinigung}}$
        — Domänenmodelle $\to$ domänenübergreifende Struktur $\to$ einheitliches Framework.
\end{itemize}

Ein lebendiges theoretisches Kontinuum erfordert:

\[
\oint_{C_{\mathrm{Kohärenz}}} dA \ungefähr 0.
\]

% ---------------------------------------------------------------
\subsubsection*{7. Grenze \texorpdfstring{$\partial\Omega(M_9)$}{\partial\Omega(M_9)}}

Bedingungen, die sich der Grenze nähern:

\begin{itemize}
    \item Aufschlüsselung des Inferenzabschlusses,
    \item Kohärenzverlust: $d_{\mathrm{coh}}>\Theta_{\mathrm{coh}}$,
    \item Verbreitung von Paradoxien,
    \item Divergenz formaler Sprachen jenseits gegenseitiger Interpretierbarkeit,
    \item Versagen der Modell-Welt-Korrespondenz (prädiktiver Kollaps),
    \item außer Kontrolle geratene theoretische Spannung $T_9$.
\end{itemize}

Das Überqueren von $\partial\Omega(M_9)$ kollabiert die Theorie zu fragmentierten $M_8$-zulässigen Strukturen.

% ---------------------------------------------------------------
\subsubsection*{8. Beziehung zu \texorpdfstring{$M_8$}{M_8} und $M_{10}$}

\textbf{Von $M_8$ bis $M_9$:}

Notwendige Voraussetzungen:

\[
A_{\mathrm{Syntax}}\neq\emptyset,
\quad
A_{\mathrm{Semantik}}\neq\emptyset,
\quad
\Theta_{\mathrm{coh}}\ \text{endlich}.
\]

Symbolische und wissenschaftliche Strukturen von $M_8$ heben sich hervor
Modellräume von $M_9$ über den Operator $\Psi$.

\textbf{In Richtung $M_{10}$ (metatheoretischer Raum):}

$M_{10}$ stellt vor:

\begin{itemize}
    \item-Kategorien von Modellkategorien,
    \item Funktionsdynamik,
    \item Metakohärenzschwellen,
    \item Zyklen der Selbstbeschreibung,
    \item der Raum, in dem Metatheorien selbst Kontinua bilden.
\end{itemize}

Somit stellt $M_9$ den strukturierten Modellraum bereit, der zum Substrat für $K_{10}$ wird.
